\section{2025-02-19}

\startframedtask
Teams: Allgemein > Datein > Ph Q1_Q2 > Semester_4_elektromagner Schwingkreis > Kernphysik > 01 und 02
\stopframedtask

\subsection{Aufgabe 1: Rutherford - Atommodel}

\startbuffer[kernphysik:atommodel:rutherford]
\startMPpage

u := 1cm;

draw fullcircle scaled 4u;

draw fullcircle scaled 1u;
label ("+", (0, 0));


draw fullcircle scaled 0.5u shifted (1u, 1u);
label ("-", (1u, 1u));

draw fullcircle scaled 0.5u shifted (-0.5u, 1u);
label ("-", (-0.5u, 1u));

draw fullcircle scaled 0.5u shifted (-0.8u, -0.8u);
label ("-", (-0.8u, -0.8u));

draw fullcircle scaled 0.5u shifted (1.2u, -0.3u);
label ("-", (1.2u, -0.3u));

draw fullcircle scaled 0.5u shifted (0.6u, -1.2u);
label ("-", (0.6u, -1.2u));

\stopMPpage
\stopbuffer

\placefigure[left][kernphysik:atommodel:rutherford]{Atommodel nach Rutherford}{
\typesetbuffer[kernphysik:atommodel:rutherford]
}

\startitemize[packed]
\item Positiv geladenere Kern
\startitemize[packed]
\item Besitzt einen großteil der Masse
\stopitemize
\item leichte Atomhülle, bestehend aus geladenen negativ geladenen Teilchen
\startitemize[packed]
\item Elektronen
\stopitemize
\stopitemize

\subsection{Aufgabe 2: Schreibweise Atomkerne}

\startformula
\startmatrix
\NC A \NR \NC Z \NR
\stopmatrix
X \hat{=}
\startmatrix
\NC \text{Massezahl} \hfill \NR
\NC \text{Ordnungszahl} \hfill \NR
\stopmatrix
\text{Elementsymbol also z. B.: }
\startmatrix
\NC 14 \NR
\NC 6 \NR
\stopmatrix
C
\stopformula

\startitemize[packed]
\item $Z$: Zahl der Protonen
\item $N$: Neutronenzahl
\item $A = Z + N$: ungefähre Masse 
\stopitemize

\startitemize[packed]
\item Isotop: Nuklide (Atomsorten) mit gleicher Ordnungszahl, aber unterschiedlicher Massezahl
\item Isomer: Chemische Verbindungen der gleichen Summenformel, aber unterschiedlicher chemischer Struktur
\item Isoton: zwei Kerne, die aus der gleichen Anzahl Neutronen, aber einer unterschiedlichen Anzahl Protonen aufgebaut sind, siehe Isoton.
\item Isoba
\stopitemize

\subsection{Aufgabe 3: Funktion einer Nebelkammer}
